\documentclass[12pt]{article}
\usepackage{amsmath,amssymb}
\usepackage{geometry}
\usepackage{enumitem}
\geometry{margin=0.8in}

\begin{document}
\begin{center}
  {\Large SAT Math: Hard-Module Challenge 2}
\end{center}

\section*{Instructions}
Complete all 5 questions in 20 minutes. A calculator, including a graphing calculator, is permitted. Select exactly one answer (A--D) per question. All choices are exact; do not round intermediate calculations. All variables are real. Use the blank space for work.

\noindent\textit{Original, unofficial challenge practice; not a complete adaptive module or a score predictor.}

\section*{Questions}
\begin{enumerate}[leftmargin=*,itemsep=1em]
  \item The quadratic function $f$ satisfies $f(2)=f(10)=15$ and $f(6)=-17$. The function $g$ is defined by $g(x)=f(2x-1)$. What is the product of the two solutions of $g(x)=0$?
  \begin{enumerate}[label=\Alph*.,itemsep=0pt]
    \item $\dfrac{55}{8}$
    \item $\dfrac{81}{8}$
    \item $\dfrac{55}{2}$
    \item $\dfrac{81}{2}$
  \end{enumerate}

  \vfill
  \item Data set A has a mean of 72, and data set B has a mean of 90. When all the values in both sets are combined, the mean is 84. The same positive number of values is then removed from each set. Both sets remain nonempty, and the mean of each individual set is unchanged. The mean of all the remaining values combined is 86. What fraction of the original total number of values was removed?
  \begin{enumerate}[label=\Alph*.,itemsep=0pt]
    \item $\dfrac15$
    \item $\dfrac13$
    \item $\dfrac35$
    \item $\dfrac25$
  \end{enumerate}

  \vfill
  \newpage
  \item At time $t=0$, culture A contains 400 bacteria and culture B contains 900 bacteria. During each hour, the number of bacteria in culture A increases by $p\%$, and the number in culture B decreases by $p\%$, where $0<p<100$. After 2 hours, the two cultures contain equal numbers of bacteria. Under this model, what is the total number of bacteria in the two cultures after 3 hours?
  \begin{enumerate}[label=\Alph*.,itemsep=0pt]
    \item $576$
    \item $1{,}152$
    \item $1{,}200$
    \item $1{,}300$
  \end{enumerate}

  \vfill
  \item Two distinct circles in the $xy$-plane each have their center in the first quadrant. Each circle is tangent to the $x$-axis, the $y$-axis, and the line $3x+4y=48$. What is the sum of the areas of the two circles?
  \begin{enumerate}[label=\Alph*.,itemsep=0pt]
    \item $16\pi$
    \item $96\pi$
    \item $592\pi$
    \item $784\pi$
  \end{enumerate}

  \vfill
  \item For which of the following ranges of $k$ does the equation
  \[
    \sqrt{x+6}=x-k
  \]
  have exactly two distinct real solutions?
  \begin{enumerate}[label=\Alph*.,itemsep=0.4em]
    \item $-\dfrac{25}{4}<k\le -6$
    \item $-\dfrac{25}{4}\le k\le -6$
    \item $-\dfrac{25}{4}<k<-6$
    \item $k>-\dfrac{25}{4}$
  \end{enumerate}
  \vfill
\end{enumerate}

\end{document}
